\chapter{Volumetric path tracing}
\label{ch:vpt}

We already introduced volumetric data and rendering in previous chapters,
but now we will go into more detail before presenting our own method in
\Cref{ch:methods}. First, we formally describe volumetric data, then we
introduce delta tracking, the sampling technique used to handle heterogeneous
media, and finally we present the volumetric path tracer that our work builds
on. For reference, \Cref{tab:notation} lists the notation used throughout this and the following chapters.

\begin{table}[p]
    \renewcommand{\arraystretch}{0.93}
    \centering
    \footnotesize
    \caption{Summary of used mathematical notation.}
    \label{tab:notation}
    \begin{tabular}{ll}
        \textbf{Symbol}                 & \textbf{Description} \\
        \midrule
        $\mathbf{x}$                    & Position vector \\
        $\omega$                        & Direction vector \\
        $v(\mathbf{x})$                 & Volumetric data value \\
        $(\omega \cdot \nabla)$         & Directional derivative \\
        $\int_{S^2}$                    & Integral over a spherical domain \\
        $\langle \cdot \rangle$         & Monte Carlo estimator \\
        $\log$                          & Natural logarithm \\
        $\xi$                           & Random number \\
        $\mathbf{x}_t$                  & Point along a ray \\
        $t$, $d$                        & Ray parameters \\
        $t \sim q(t)$                   
            & $t$ distributed according to probability density $q$ \\
        \hline        
        $L(\mathbf{x}, \omega)$         & Radiance field \\
        $L_s(\mathbf{x}, \omega)$       & In-scattering radiance field \\
        $L_e(\mathbf{x}, \omega)$       & Emission radiance field \\  
        $T(t)$                          & Transmittance \\

        \hline         
        $\sigma_a(\mathbf{x})$          & Absorption coefficient \\ 
        $\sigma_s(\mathbf{x})$          & Scattering coefficient \\ 
        $\sigma_t(\mathbf{x}) = \sigma_a(\mathbf{x}) + \sigma_s(\mathbf{x})$
                                        & Extinction coefficient \\ 
        $\sigma_n(\mathbf{x})$          & Null-collision coefficient \\ 
        $\overline\sigma = \sigma_t(\mathbf{x}) + \sigma_n(\mathbf{x})$
                                        & Free-path coefficient \\
        \hline
        $P_a(\mathbf{x})$               & Absorption probability \\
        $P_s(\mathbf{x})$               & Scattering probability \\
        $P_n(\mathbf{x})$               & Null-collision probability \\
        \hline
        $f_p(\mathbf{x}, \omega, \omega')$
                                        & Phase function \\
        $p(\theta)$                     & Henyey-Greenstein phase function \\
        $g$                             & Anisotropy or asymmetry parameter \\ 
        $\theta(\omega)$                & Angle of vector $\omega$ \\
        $f_t(v)$                        & Transfer function \\
        $(\cdot)^{(a)}, (\cdot)^{(r)}, (\cdot)^{(g)}, (\cdot)^{(b)}$
            & Alpha, red, green and blue component of a color \\
        \hline
        $L$                             & Number of encoding levels \\
        $F$                             & Encoding feature dimension \\
        $N_{\min}^{\text{pos}}$, $N_{\max}^{\text{pos}}$,
        $N_{\min}^{\text{dir}}$, $N_{\max}^{\text{dir}}$
                                        & Resolutions for
                                          position and direction encodings \\
        $T$                             & Maximum encoding table size \\
        \hline
        $\mathbf{c}$, $\mathbf{c}_n$    & RGB color triplet \\
        $\sigma_k$, $\sigma$, $\tau$    & Bilateral filter parameters \\
        $r$                             & Bilateral filter radius \\
        $w_d$, $w_r$                    & Spatial and range weights \\
    \end{tabular}
\end{table}

\section{Volumetric data}

A set of samples $(x, y, z, v)$ in 3D space is called \emph{volumetric data}.
These samples, also called \emph{voxels} when positioned on a regular grid,
represent a value $v$ at the 3D location $(x, y, z)$. The data is binary if
$v$ is either $0$ or $1$, but values can also be integers from a set,
multivalued, or even vectors. Values usually represent some measurable property,
such as color or velocity, as in \ac{mri} and \ac{ct}
imaging~\cite{kaufman_7_2005, kaufman_43_2000}.

Samples can be taken at random locations, but if they are taken at constant 
intervals along all orthogonal axes, the data is \emph{isotropic}. If each
axis has a different spacing, the data is \emph{anisotropic} instead. An example
of anisotropic volumetric data is medical imaging, because image slices are
usually spaced differently from in-plane samples. Both isotropic and
anisotropic sets are defined on a regular grid and can be stored in a 3D array.
We call such sets simply \emph{volumes}. However, not all medical imaging
modalities are isotropic. An example of such a modality is 3D ultrasound,
which is defined on a non-Cartesian grid, such as the acoustic
grid~\cite{fernando_chapter_2004}.

\section{Light in volumes}

When rendering an image, our goal is to compute the amount of light entering
the camera when looking in a certain direction. In surface rendering, we send
rays into the scene and bounce them off the surfaces. When working with volumes,
this is not enough because light also interacts with the particles inside
the volume itself. When a photon hits a volume, it begins interacting with
particles inside the volume. These collisions define the volume radiance
distribution~\cite{fong_production_2017, wrenninge_production_2012}.

\begin{figure}[b]
    \centering
    \begin{subfigure}{0.23\textwidth}
    \centering
    \begin{tikzpicture}
        \node at (0,1.1) {\scriptsize Absorption}; 
        
        \filldraw[draw=cyan!15, fill=cyan!5, thick] (0,0) circle (0.6cm);
        \fill (-0.1,0.3) circle (1pt); 
        \fill (0.3,0.15) circle (1pt); 
        \fill (0.05,-0.2) circle (1pt); 
        \draw[-, thin] (-1,0.3) -- (-0.1,0.3);
        \draw[-, thin] (-1.1,0.15) -- (0.3,0.15);
        \draw[-, thin] (-0.9,-0.2) -- (0.05,-0.2);
        \draw[{Stealth[length=4pt]}-{Stealth[length=4pt]}, semithick]
            (-1.3,0) node[below] {\scriptsize $\omega$} --
            (1.3,0) node[below] {\scriptsize $L$};
    \end{tikzpicture}
\end{subfigure}%
\begin{subfigure}{0.23\textwidth}
    \centering
    \begin{tikzpicture}
        \node at (0,1.1) {\scriptsize Out-scattering}; 
        
        \filldraw[draw=cyan!15, fill=cyan!5, thick] (0,0) circle (0.6cm);
        \fill (-0.1,0.3) circle (1pt); 
        \fill (0.3,0.15) circle (1pt); 
        \fill (0.05,-0.2) circle (1pt); 
        \draw[-, thin] (-1,0.3) -- (-0.1,0.3);
        \draw[-, thin] (-1.1,0.15) -- (0.3,0.15);
        \draw[-, thin] (-0.9,-0.2) -- (0.05,-0.2);
        \draw[-{Stealth[length=3pt]}, thin] (-0.1,0.3) -- (0.5,0.7);
        \draw[-{Stealth[length=3pt]}, thin] (0.3,0.15) -- (0.8,0.4);
        \draw[-{Stealth[length=3pt]}, thin] (0.05,-0.2) -- (0.75,-0.45);
        \draw[{Stealth[length=4pt]}-{Stealth[length=4pt]}, semithick]
            (-1.3,0) node[below] {\scriptsize $\omega$} --
            (1.3,0) node[below] {\scriptsize $L$};
    \end{tikzpicture}
\end{subfigure}%
\begin{subfigure}{0.23\textwidth}
    \centering
    \begin{tikzpicture}
        \node at (0,1.1) {\scriptsize Emission}; 
        
        \filldraw[draw=cyan!15, fill=cyan!5, thick] (0,0) circle (0.6cm);
        \fill (-0.1,0.3) circle (1pt); 
        \fill (0.3,0.15) circle (1pt); 
        \fill (0.05,-0.2) circle (1pt); 
        \draw[-{Stealth[length=3pt]}, thin] (-0.1,0.3) -- (0.9,0.3);
        \draw[-{Stealth[length=3pt]}, thin] (0.3,0.15) -- (1,0.15);
        \draw[-{Stealth[length=3pt]}, thin] (0.05,-0.2) -- (0.8,-0.2);
        \draw[{Stealth[length=4pt]}-{Stealth[length=4pt]}, semithick]
            (-1.3,0) node[below] {\scriptsize $\omega$} --
            (1.3,0) node[below] {\scriptsize $L$};           
    \end{tikzpicture}
\end{subfigure}%
\begin{subfigure}{0.23\textwidth}
    \centering
    \begin{tikzpicture}
        \node at (0,1.1) {\scriptsize In-scattering}; 
        
        \filldraw[draw=cyan!15, fill=cyan!5, thick] (0,0) circle (0.6cm);
        \fill (-0.1,0.3) circle (1pt); 
        \fill (0.3,0.15) circle (1pt); 
        \fill (0.05,-0.2) circle (1pt); 
        \draw[-, thin] (-0.8,0.6) -- (-0.1,0.3);
        \draw[-, thin] (-0.9,0.2) -- (0.3,0.15);
        \draw[-, thin] (-0.7,-0.4) -- (0.05,-0.2);
        \draw[-{Stealth[length=3pt]}, thin] (-0.1,0.3) -- (0.9,0.3);
        \draw[-{Stealth[length=3pt]}, thin] (0.3,0.15) -- (1,0.15);
        \draw[-{Stealth[length=3pt]}, thin] (0.05,-0.2) -- (0.8,-0.2);
        \draw[{Stealth[length=4pt]}-{Stealth[length=4pt]}, semithick]
            (-1.3,0) node[below] {\scriptsize $\omega$} --
            (1.3,0) node[below] {\scriptsize $L$};
    \end{tikzpicture}
\end{subfigure}%
    \caption{Four possible light interaction events inside volumes.}
    \label{fig:rte}
\end{figure}

\begin{figure}[b]
    \centering
    \begin{subfigure}{0.3\textwidth}
        \centering
        \includegraphics[width=0.85\textwidth]{absorption.png}
    \end{subfigure}
    \begin{subfigure}{0.3\textwidth}
        \centering
        \includegraphics[width=0.85\textwidth]{scattering.png}
    \end{subfigure}
    \begin{subfigure}{0.3\textwidth}
        \centering
        \includegraphics[width=0.85\textwidth]{emission.png}
    \end{subfigure}
    
    \caption{
        Examples of a purely absorptive (left), scattering (middle) and
        emissive (right) volume~\cite{wrenninge_production_2011}.
    }
    \label{fig:interactions}
\end{figure}

More formally, we want to obtain radiance $L$ when looking in
direction~$\omega$. This is done using the \emph{\ac{rte}}, which defines the
radiance distribution and represents the equilibrium radiance field $L$ as a
function of two parameters, position $\mathbf{x}$ and direction $\omega$. It
contains four terms, which are visualized in \Cref{fig:rte}. Each term
contributes to the change in radiance, which can be described using a
directional derivative:
\begin{itemize}
    \item \emph{Absorption} is the loss of radiance when a photon is absorbed by
    the medium. It is proportional to the radiance at that point, and using the
    \emph{absorption coefficient} $\sigma_a$ we write:
    \begin{equation}
        (\omega \cdot \nabla) L(\mathbf{x}, \omega) =
            -\sigma_a(\mathbf{x}) L(\mathbf{x}, \omega).
    \end{equation}

    \item \emph{Out-scattering} is the loss of radiance when a photon hits a
    particle that reflects it in a direction other than $\omega$. It is also
    proportional to the radiance at that point, but here we use the
    \emph{scattering coefficient} $\sigma_s$:
    \begin{equation}
        (\omega \cdot \nabla) L(\mathbf{x}, \omega) =
            -\sigma_s(\mathbf{x}) L(\mathbf{x}, \omega).
    \end{equation}

    \item \emph{Emission} occurs when molecules are hot enough to emit photons
    in all directions. If $L_e$ is the radiance field of emitted light, we
    write:
    % The radiance that is emitted is a separate radiance field $L_e$ using which we write:
    \begin{equation}
        (\omega \cdot \nabla) L(\mathbf{x}, \omega) =
            \sigma_a(\mathbf{x}) L_e(\mathbf{x}, \omega). 
    \end{equation}
    Depending on the derivation, some authors may omit the absorption
    coefficient here.

    \item \emph{In-scattering} is the reverse of out-scattering. It is a
    contribution of radiance when a photon traveling in some other direction
    $\omega'$ hits a particle and is reflected into the direction $\omega$.
    The angular distribution of scattering, i.e., the probability of scattering 
    from angle $\omega'$ to $\omega$, is modeled by a
    \emph{phase function} $f_p(\mathbf{x}, \omega, \omega')$. Because we have
    to take into account all the angles from which photons can travel,
    we write this term using an integral over the entire spherical domain:
    \begin{align}
        (\omega \cdot \nabla) L(\mathbf{x}, \omega) &=
            \sigma_s(\mathbf{x})
            \int_{\mathit{S}^2} f_p(\mathbf{x}, \omega, \omega')
            L(\mathbf{x}, \omega')\, \mathrm{d}\omega' \\
        &= \sigma_s(\mathbf{x}) L_s(\mathbf{x}, \omega).
    \end{align}
    An additional term $L_s$ is also introduced to make the notation more concise.
\end{itemize}
In \Cref{fig:interactions}, we see how each of these interactions by itself
affects the same volume.

The final \ac{rte} is obtained by combining all of the terms:
\begin{equation}
    (\omega \cdot \nabla) L(\mathbf{x}, \omega) =
        -\sigma_t(\mathbf{x}) L(\mathbf{x}, \omega)
        +\sigma_a(\mathbf{x}) L_e(\mathbf{x}, \omega) 
        +\sigma_s(\mathbf{x}) L_s(\mathbf{x}, \omega),
\end{equation}
where $\sigma_t = \sigma_a + \sigma_s$ is the \emph{extinction coefficient},
This substitution is made because of the similarity between absorption and
out-scattering.

Before writing the final rendering equation, we introduce \emph{transmittance}
\begin{equation}
    \label{eq:transmittance}
    T(t) = \exp \left( -\int_0^t \sigma_t(\mathbf{x}_s)\, \mathrm{d}s \right),
\end{equation}
which is the reduction factor between $\mathbf{x}$ and a position on the ray 
$\mathbf{x}_t = \mathbf{x} - t\omega$ due to absorption and out-scattering.
We use it after integrating the \ac{rte} on both sides to write the
\emph{\ac{vre}}:
\begin{align}
    \label{eq:vre}
    L&(\mathbf{x}, \omega) =
        \int_0^d T(t) \bigg[
            \sigma_a(\mathbf{x}_t) L_e(\mathbf{x}_t, \omega)  
            + \sigma_s(\mathbf{x}_t) L_s(\mathbf{x}_t, \omega)
            + L(\mathbf{x}_d, \omega)
        \bigg]\, \mathrm{d}t \\
    &= \int_0^d T(t) \bigg[
            \sigma_a(\mathbf{x}_t) L_e(\mathbf{x}_t, \omega)  
            + \sigma_s(\mathbf{x}_t) L_s(\mathbf{x}_t, \omega)
        \bigg]\, \mathrm{d}t
        + T(d) L(\mathbf{x}_d, \omega).
\end{align}
The integrating ray ends at parameter $d$, where it hits the surface or another
volume.

\section{Delta tracking}

If we focus on scenes without surfaces, we only have to solve the integral in
the \ac{vre}. For this, we use Monte Carlo integration. We can write a
single-point Monte Carlo estimator that we need to
evaluate~\cite{novak_monte_2018}:
\begin{equation}
    \left\langle L(\mathbf{x}, \omega) \right\rangle =
        \frac{T(t)}{q(t)} 
        \bigg[
            \sigma_a(\mathbf{x}_t) L_e(\mathbf{x}_t, \omega)  
            + \sigma_s(\mathbf{x}_t) L_s(\mathbf{x}_t, \omega)
        \bigg],
\end{equation}
where $t$ is sampled according to distribution $q$. This distribution should
be as close to $T$ as possible to reduce the estimator variance. In
homogeneous volumes with constant $\sigma_t$, transmittance can be simplified,
inverted, and importance sampled. However, in general, volumes have 
heterogeneous density, and we have to use \emph{delta tracking}, sometimes
called \emph{Woodcock tracking}~\cite{woodcock_techniques_1965}.

The idea is to introduce a fictitious volume that, together with our original
volume, forms a homogeneous volume. We choose a \emph{null-collision coefficient}
$\sigma_n$ such that the \emph{free-path coefficient}
\begin{equation}
    \overline\sigma
        = \sigma_a(\mathbf{x}) + \sigma_s(\mathbf{x}) + \sigma_n(\mathbf{x})
        = \sigma_t(\mathbf{x}) + \sigma_n(\mathbf{x}) 
\end{equation}
is constant. It follows that the transmittance in \Cref{eq:transmittance}
simplifies to
\begin{equation}
    T_{\overline\sigma}(t) = \exp(-t \overline\sigma).
\end{equation}
\Cref{eq:vre} can now be written as
\begin{align}
\begin{split}
    \label{eq:null_vre}
    L(\mathbf{x}, \omega) =
        \int_0^d T_{\overline\sigma}(t) \bigg[
            &\sigma_a(\mathbf{x}_t) L_e(\mathbf{x}_t, \omega)
            + \sigma_s(\mathbf{x}_t) L_s(\mathbf{x}_t, \omega) \\
            &+ \sigma_n(\mathbf{x}_t) L(\mathbf{x}_t, \omega)
            + L(\mathbf{x}_d, \omega)
        \bigg]\, \mathrm{d}t.
\end{split}
\end{align}
Alongside the existing terms for emission or absorption and in-scattering,
we now have an additional term for null collisions. The estimator for this
equation without the surface term is now
\begin{align}
\begin{split}
    \label{eq:vre_mc}
    \left\langle L(\mathbf{x}, \omega) \right\rangle =
        \frac{T_{\overline\sigma}(t)}{q_{\overline\sigma}(t)} 
        \bigg[
            \sigma_a(\mathbf{x}_t) L_e(\mathbf{x}_t, \omega)  
            &+ \sigma_s(\mathbf{x}_t) L_s(\mathbf{x}_t, \omega) \\
            &+ \sigma_n(\mathbf{x}_t) L(\mathbf{x}_t, \omega)
        \bigg].
\end{split}
\end{align}
The total extinction coefficient is constant, so we know how to compute
$T_{\overline\sigma}$ and sample proportionally to it:
\begin{equation}
    \label{eq:q_pdf}
    t \sim q_{\overline\sigma}(t) = \overline\sigma \exp(-t \overline\sigma).
\end{equation}
With inverse sampling, we can sample a random number $\xi \in (0, 1]$ and compute
\begin{equation}
    t(\xi) = -\frac{\log(\xi)}{\overline\sigma}.
\end{equation}

After sampling a distance $t$ of a possible collision event, there are three
components in \Cref{eq:vre_mc} to evaluate. In delta tracking, we classify
the event as either a null collision, scattering, or absorption based on the
probabilities we get if we express probability density in terms of transmittance
and insert it into \Cref{eq:vre_mc}:
\begin{align}
    \label{eq:delta}
    \left\langle L(\mathbf{x}, \omega) \right\rangle &=
        \frac{\sigma_a}{\overline\sigma} L_e(\mathbf{x}_t, \omega)  
        + \frac{\sigma_s}{\overline\sigma} L_s(\mathbf{x}_t, \omega)
        + \frac{\sigma_n}{\overline\sigma} L(\mathbf{x}_t, \omega) \\
    &= P_a(\mathbf{x}_t) L_e(\mathbf{x}_t, \omega)  
        + P_s(\mathbf{x}_t) L_s(\mathbf{x}_t, \omega)
        + P_n(\mathbf{x}_t) L(\mathbf{x}_t, \omega).
\end{align}
The probabilities are thus based on ratios between coefficients:
\begin{equation}
    \label{eq:probs}
    P_a(\mathbf{x}) = \frac{\sigma_a(\mathbf{x})}{\overline\sigma}, \qquad
    P_s(\mathbf{x}) = \frac{\sigma_s(\mathbf{x})}{\overline\sigma}, \qquad
    P_n(\mathbf{x}) = \frac{\sigma_n(\mathbf{x})}{\overline\sigma}.
\end{equation}

\begin{figure}[t]
    \centering
    \begin{tikzpicture}[scale=1.25]
    % --- fills (painted back to front: gray, then blue, then red) ---
    % top region: whole box in faded gray
    \fill[gray!15] (0,0) rectangle (8,3);
    % middle region: everything below the upper curve in faded blue
    \fill[cyan!20] (0,2.55)
        .. controls (0.4,2.0)  and (0.7,2.04)  .. (1.04,2.02)
        .. controls (1.45,2.00) and (1.69,1.88) .. (2.09,1.52)
        .. controls (2.54,1.09) and (3.34,0.52) .. (4.49,1.09)
        .. controls (5.2,1.4)   and (5.2,2.20)  .. (5.61,2.22)
        .. controls (6.1,2.23)  and (6.15,1.35) .. (6.75,1.37)
        .. controls (7.4,1.39)  and (7.38,1.9)  .. (8.0,2.02)
        -- (8,0) -- (0,0) -- cycle;
    % bottom region: everything below the lower curve in faded red
    \fill[red!20] (0,0.77)
        .. controls (0.67,1.26) and (1.23,1.34) .. (1.92,1.07)
        .. controls (2.74,0.76) and (3.30,0.17) .. (4.28,0.60)
        .. controls (4.64,0.76) and (5.03,0.83) .. (5.43,0.77)
        .. controls (5.91,0.70) and (6.24,0.34) .. (6.84,0.36)
        .. controls (7.41,0.39) and (7.67,0.76) .. (8.0,1.22)
        -- (8,0) -- (0,0) -- cycle;

    % --- outlines ---
    \draw[thick, cyan!60] (0,2.55)
        .. controls (0.4,2.0)  and (0.7,2.04)  .. (1.04,2.02)
        .. controls (1.45,2.00) and (1.69,1.88) .. (2.09,1.52)
        .. controls (2.54,1.09) and (3.34,0.52) .. (4.49,1.09)
        .. controls (5.2,1.4)   and (5.2,2.20)  .. (5.61,2.22)
        .. controls (6.1,2.23)  and (6.15,1.35) .. (6.75,1.37)
        .. controls (7.4,1.39)  and (7.38,1.9)  .. (8.0,2.02);
    \draw[thick, red!60] (0,0.77)
        .. controls (0.67,1.26) and (1.23,1.34) .. (1.92,1.07)
        .. controls (2.74,0.76) and (3.30,0.17) .. (4.28,0.60)
        .. controls (4.64,0.76) and (5.03,0.83) .. (5.43,0.77)
        .. controls (5.91,0.70) and (6.24,0.34) .. (6.84,0.36)
        .. controls (7.41,0.39) and (7.67,0.76) .. (8.0,1.22);
    \draw[thick] (0,3) -- (8,3) -- (8,0);

    % --- text ---
    \node[gray!80] at (7.78,2.5) {\footnotesize $\sigma_n$}; 
    \node[cyan!80] at (7.78,1.5) {\footnotesize $\sigma_s$}; 
    \node[red!80] at (7.78,0.5) {\footnotesize $\sigma_a$}; 
    \node at (-0.2,3)   {\footnotesize $\overline\sigma$}; 

    % --- sampling ---
    % named points
    \coordinate (x0) at (0,   1.35);  % start of x -> x_t arrow
    \coordinate (x1) at (1.8, 2.2);   
    \coordinate (x2) at (3.4, 0.7);
    \coordinate (x3) at (5.0, 2.0);
    \coordinate (x4) at (6.6, 0.2);
    \def\lbly{-0.25}                   % y of the bottom-axis labels
    
    \draw[-Stealth, semithick] (x0) -- ([xshift=-0.05cm] x1 |- x0)
        node[midway, above] {\footnotesize $t \sim \mathrm{Exp}(\overline\sigma)$};
    \draw[Stealth-Stealth, semithick] (x1 |- 0,0.05) -- (x1 |- 0,2.95);
    
    \fill (x1) circle[radius=1.5pt];
    \fill (x2) circle[radius=1.5pt];
    \fill (x3) circle[radius=1.5pt];
    \fill (x4) circle[radius=1.5pt];
    
    \node at (x0 |- 0,\lbly) {\footnotesize $\mathbf{x}_0$};
    \node at (x1 |- 0,\lbly) {\footnotesize $\mathbf{x}_1$};
    \node at (x2 |- 0,\lbly) {\footnotesize $\mathbf{x}_2$};
    \node at (x3 |- 0,\lbly) {\footnotesize $\mathbf{x}_3$};
    \node at (x4 |- 0,\lbly) {\footnotesize $\mathbf{x}_4$};
    
    % --- axes ---
    \draw[Stealth-Stealth, very thick] (0,3.5) -- (0,0) -- (8.5,0);
\end{tikzpicture}
    \caption{
        Four steps of delta tracking in a heterogeneous volume. At each step,
        the step size is sampled from an exponential distribution and a D3 die
        is tossed. We either hit the fictitious volume ($\mathbf{x}_1$,
        $\mathbf{x}_3$) or the real volume with a scattering ($\mathbf{x}_2$)
        or an absorption event ($\mathbf{x}_4$).
    }
    \label{fig:delta}
\end{figure}

Evaluating \Cref{eq:delta} exactly would require computing all three terms
on the right-hand side. Instead, we can use Russian roulette for each term,
evaluating it and dividing by $P$ to keep the estimator unbiased with
probability $P$ from \Cref{eq:probs}, and treating it as $0$ with probability
$1 - P$. Applied to each term, this would still require evaluating each term
independently. However, since
$P_a(\mathbf{x}) + P_s(\mathbf{x}) + P_n(\mathbf{x}) = 1$, we can correlate
all three Russian roulettes into a single random decision and evaluate only
one term without affecting the result.

In practice, this means that at each ray-tracing step, we roll a D3 die with
these probabilities. If a null collision is selected, we hit the fictitious
volume, so we move forward without changing direction. Otherwise, we hit the
actual volume. In the case of scattering, we still move forward but also change
direction according to the phase function. In the case of absorption, ray 
tracing is stopped. Four steps of this process are visualized in
\Cref{fig:delta}.

\section{VPT}

We base our work on an existing volumetric path tracer,
\emph{VPT}~\cite{lesar_real-time_2018}. It runs in a browser and is therefore
cross-platform. The initial implementation used the WebGL API for graphics
programming, but we use the updated version rewritten for the newer WebGPU API.

It implements delta tracking with multiple scattering and can be configured using:
\begin{itemize}
    \item \emph{Extinction.} Rate of the exponential distribution used when
    sampling the step size $t$.

    \item \emph{Anisotropy.} Parameter $g \in (-1, 1)$ of the Henyey-Greenstein
    phase function, which models the scattering angle:
    \begin{equation}
        p(\theta) = \frac{1}{4\pi}
            \frac{1 - g^2}{(1 + g^2 - 2g\cos\theta)^{3/2}}.
    \end{equation}
    Depending on the value of $g$, we either have isotropic scattering ($g = 0$),
    forward scattering ($g > 0$), or backward scattering
    ($g < 0$)~\cite{fong_production_2017}.

    \item \emph{Bounces.} Upper bound on the number of real scattering events
    in a single iteration.
    
    \item \emph{Steps.} Upper bound on the number of events (null collision,
    scattering, or absorption) in a single iteration.

    \item \emph{Transfer function.} Mapping $f_t$ of grayscale volume values
    $v \in [0, 255]$ to color. This allows us to modify the overall color of
    the volume, either on its surface or inside it.
\end{itemize}

The path tracer determines the pixel color by tracing a ray until it goes out
of bounds. In the case of an absorption event, the pixel is colored black.
During scattering events, the color contribution $f_t(v)$ is accumulated into
the transmittance, and the direction $\omega$ is updated according to
$p(\theta(\omega))$. Null collision events do not change the ray’s path, and
if the ray reaches the boundary of the scene, the pixel color is computed from
the accumulated transmittance. At each interaction point, the type of event is
chosen stochastically according to the following probabilities:
\begin{align}
    P_n(\mathbf{x}) &= 1 - f_t(v(\mathbf{x}))^{(a)}, \\
    P_s(\mathbf{x}) &= f_t(v(\mathbf{x}))^{(a)} \cdot \max\{
        f_t(v(\mathbf{x}))^{(r)},
        f_t(v(\mathbf{x}))^{(g)},
        f_t(v(\mathbf{x}))^{(b)}
    \}, \\
    P_a(\mathbf{x}) &= 1 - P_n(\mathbf{x}) - P_s(\mathbf{x}).
\end{align}
Here, the superscripts $(r)$, $(g)$, $(b)$, and $(a)$ denote the red, green,
blue, and alpha channels of the RGBA color returned by the transfer function
$f_t$, respectively.